\part{Quantum Information and Noncommutative Analysis}

\chapter{Composite tests require matrix norms}
\label{ch:operator-spaces}

\physicalquestion{Why can two linear maps with the same ordinary norm behave
differently when an entangled auxiliary system is attached?}

\modelassumptions{Systems are finite dimensional in the explicit
calculations.  A subspace of operators inherits the norm arising from its
action on a Hilbert space, and an $r$-level ancilla is represented by matrix
amplification.}

\section{Which norm remembers every ancilla?}

A concrete operator space is a closed subspace
$E\subseteq\mathcal B(H,K)$.  Its $r$th matrix level is
\[
M_r(E)\subseteq\mathcal B(H^{\oplus r},K^{\oplus r}).
\]
For $\phi:E\to F$, set $\phi_r([x_{ij}])=[\phi(x_{ij})]$ and
\[
\|\phi\|_{\mathrm{cb}}=\sup_r\|\phi_r\|.
\]
The row and column operator spaces are
$R_n=M_{1,n}$ and $C_n=M_{n,1}$ with their concrete matrix norms.  As Hilbert
spaces both are $\C^n$; as operator spaces they encode different directions
of action.

\section{Can ordinary norm miss an arbitrarily large effect?}

\begin{theorem}[Transposition detects matrix levels]
For transposition $t:M_n\to M_n$,
\[
\|t\|=1,
\qquad
\|t\|_{\mathrm{cb}}=n.
\]
Thus ordinary boundedness does not control behavior under ancillas.
\end{theorem}

\begin{proof}
Transposition preserves singular values, so its operator norm on $M_n$ is
one.  At matrix level $n$, apply $\id\otimes t$ to the flip
$F=\sum_{ij}E_{ij}\otimes E_{ji}$, which is unitary and has norm one.  The
image is
$\sum_{ij}E_{ij}\otimes E_{ij}=n|\Omega\rangle\langle\Omega|$ for normalized
$\Omega=n^{-1/2}\sum_i e_i\otimes e_i$, and has norm $n$.  Hence the cb norm
is at least $n$.  Writing a matrix amplification as block transposition and
using the row--column factorization gives the reverse bound $n$.
\end{proof}

\section{What does this mean for composite quantum systems?}

The transpose preserves positivity of an isolated density matrix but, as
seen in Chapter~\ref{ch:channels}, fails positivity on an entangled ancilla.
The norm calculation measures the same obstruction quantitatively: the gap
between bounded and completely bounded behavior grows with dimension.  Matrix
levels are therefore operational data, not decorative copies of a Banach
norm.

\section{What further contact should be proved?}

\begin{exercise}[Rows and columns are completely different]
Prove that the identity maps between the underlying Hilbert spaces of $R_n$
and $C_n$ have cb norm $\sqrt n$ in both directions.
\end{exercise}

\begin{exercise}[Ruan's concrete identities]
For $x\in M_m(E)$ and $y\in M_n(E)$ prove
\[
\|x\oplus y\|=\max\{\|x\|,\|y\|\},\qquad
\|\alpha x\beta\|\leq\|\alpha\|\|x\|\|\beta\|.
\]
Explain why these identities are invariant under complete isometry.
\end{exercise}


\chapter{Noncommutative products require Haagerup tensoring}
\label{ch:haagerup}

\physicalquestion{Which tensor product records the norm of a process formed
by multiplying or composing noncommuting operator-valued inputs?}

\modelassumptions{Operator spaces are concrete and finite dimensional for
examples, although the definitions extend by completion.  Bilinear maps are
tested at all matrix levels using matrix multiplication of indices.}

\section{Why does the algebraic tensor product forget too much?}

For $u=\sum_{k=1}^r x_k\otimes y_k\in E\otimes F$, define
\[
\|u\|_h=\inf
\left\{
\|[x_1\ \cdots\ x_r]\|_{M_{1,r}(E)}
\|[y_1\ \cdots\ y_r]^{\mathsf T}\|_{M_{r,1}(F)}
\right\}.
\]
The completion is the Haagerup tensor product $E\otimes_hF$.

For a bilinear map $b:E\times F\to G$, define
\[
b_n(x,y)_{ij}=\sum_kb(x_{ik},y_{kj}),
\qquad \|b\|_{\mathrm{mb}}=\sup_n\|b_n\|.
\]

\section{Which universal property fits multiplication?}

\begin{theorem}[Haagerup linearization]
A bilinear map $b:E\times F\to G$ is multiplicatively bounded exactly when
its linearization extends to a completely bounded map
\[
\widetilde b:E\otimes_hF\to G.
\]
Moreover $\|\widetilde b\|_{\mathrm{cb}}=\|b\|_{\mathrm{mb}}$.
\end{theorem}

\begin{proof}
A factorization $u=a\odot c$ with
$a\in M_{1,r}(E)$ and $c\in M_{r,1}(F)$ is sent at matrix level to the
corresponding product $b_n(a,c)$.  Hence
$\|\widetilde b_n(u)\|\leq\|b\|_{\mathrm{mb}}\|a\|\|c\|$; taking the
infimum proves one inequality and extension by continuity.  Conversely apply
$\widetilde b_n$ to elementary matrix factorizations $x\odot y$ to obtain
$\|b_n(x,y)\|\leq\|\widetilde b\|_{\mathrm{cb}}\|x\|\|y\|$.  Taking
suprema gives equality.
\end{proof}

\section{Where does this tensor norm appear physically?}

Multiplication in a $C^*$-algebra is multiplicatively bounded and therefore
linearizes through $\otimes_h$.  The complete isometries
\[
C_n\otimes_hR_n\cong M_n,
\qquad
R_n\otimes_hC_n\cong S_1^n
\]
distinguish operator norm from trace norm solely by reversing factors.  The
second norm controls optimal state discrimination, so the order-sensitive
tensor geometry has an operational meaning.

\section{What further contact should be proved?}

\begin{exercise}[Helstrom discrimination]
For states $\rho_0,\rho_1$ with equal prior probability, prove that the best
binary measurement succeeds with probability
\[
\frac12+\frac14\|\rho_0-\rho_1\|_1.
\]
Relate the trace norm to $R_n\otimes_hC_n$.
\end{exercise}

\begin{exercise}[Forgetting matrix norms loses tensor information]
Show that $R_n$ and $C_n$ have isometric underlying Banach spaces while their
Haagerup products in the two orders have different norms.  Explain why no
symmetric completely isometric flip can identify the two products.
\end{exercise}


\chapter{Channels are distinguished at the matrix level}
\label{ch:diamond}

\physicalquestion{How well can one use distinguish two quantum processes when
the input may be entangled with an auxiliary system?}

\modelassumptions{Channels act between finite matrix algebras, are used once,
and have equal prior probability.  Arbitrary finite-dimensional ancillas and
joint output measurements are allowed.}

\section{Which distance includes every composite probe?}

Recall complete positivity and channels from Chapter~\ref{ch:channels}.  For
a linear map $\Psi:M_n\to M_m$, its diamond norm is
\[
\|\Psi\|_\diamond
=\sup_r\|\Psi\otimes\id_{M_r}\|_{1\to1}.
\]
The supremum explicitly asks how distinguishability changes when an untouched
ancilla is attached.  Operator systems---self-adjoint unital subspaces of
$\mathcal B(H)$---supply the matrix order that makes complete positivity
intrinsic even when the domain is not a full algebra.

\section{Why is the diamond norm operational?}

\begin{theorem}[One-use channel discrimination]
For equally likely channels $\Phi_0,\Phi_1:M_n\to M_m$, the optimal one-use
success probability is
\[
p_{\mathrm{succ}}
=\frac12+\frac14\|\Phi_0-\Phi_1\|_\diamond.
\]
The ancillary dimension $n$ is sufficient.
\end{theorem}

\begin{proof}
For an input state $\rho$ on system plus ancilla, the two output states differ
by $((\Phi_0-\Phi_1)\otimes\id)(\rho)$.  The Helstrom formula from
Chapter~\ref{ch:haagerup} gives the optimal measurement for that input.
Maximizing over inputs produces the diamond norm.  Convexity reduces the
maximum to pure inputs, and Schmidt decomposition from Chapter~\ref{ch:tensors}
shows that a pure input needs Schmidt rank at most $n$; a larger ancilla cannot
improve the value.
\end{proof}

\section{What advantage can entanglement provide?}

If the maximizing input is entangled, a strategy restricted to uncorrelated
probes attains a smaller norm.  The advantage is not additional channel use;
it comes from evaluating the same map at a higher matrix level.  Process
discrimination experiments compare observed success frequencies with the
separable benchmark and the diamond-norm prediction under calibrated noise
assumptions.

\section{What further contact should be proved?}

\begin{exercise}[Two unitary channels]
For unitary channels $\rho\mapsto U_i\rho U_i^*$, express perfect one-use
distinguishability in terms of the numerical range of $U_0^*U_1$.
\end{exercise}

\begin{exercise}[The Heisenberg adjoint has cb norm]
Prove in finite dimensions that
$\|\Psi\|_\diamond=\|\Psi^*\|_{\mathrm{cb}}$ when the adjoint is equipped
with operator norm matrix levels.
\end{exercise}


\chapter{Bell correlations are noncommutative geometry}
\label{ch:bell}

\physicalquestion{How strongly can separated quantum measurements correlate
without allowing communication between the laboratories?}

\modelassumptions{Each party chooses one of two binary observables.  The
observables of different parties commute in the tensor-product model, each
outcome is $\pm1$, and the experimental settings are statistically independent
of the prepared state.}

\section{Which operator encodes the Bell test?}

For self-adjoint unitaries $A_0,A_1$ and $B_0,B_1$, define the CHSH operator
\[
S=A_0\otimes(B_0+B_1)+A_1\otimes(B_0-B_1).
\]
Local hidden-variable assignments satisfy $|\langle S\rangle|\leq2$.  The
operators need not commute within one laboratory, and that noncommutativity
changes the bound.

\section{How large can the quantum value be?}

\begin{theorem}[Tsirelson bound]
Every quantum realization above satisfies
$\|S\|\leq2\sqrt2$, and the bound is attained by a maximally entangled pair
of qubits with suitable Pauli measurements \cite{tsirelson}.
\end{theorem}

\begin{proof}
Using $A_i^2=B_j^2=1$, expand to obtain
\[
S^2=4\,1-[A_0,A_1]\otimes[B_0,B_1].
\]
Each commutator has norm at most two, so $\|S^2\|\leq8$ and
$\|S\|\leq2\sqrt2$.  For the singlet state, take orthogonal Pauli directions
for Alice and the normalized sum and difference directions for Bob; direct
calculation gives magnitude $2\sqrt2$.
\end{proof}

\section{What does a violation establish experimentally?}

Under the stated locality and statistical assumptions, an observed value
above two rules out the corresponding local hidden-variable model.  It does
not permit signaling, since each local marginal is independent of the remote
choice.  Grothendieck-type inequalities generalize the theorem: for broad
families of correlation experiments, quantum advantage over classical vector
assignments remains bounded by a universal constant.  Loophole-controlled
Bell experiments compare measured correlations with these bounds under
explicitly tested locality assumptions \cite{hensen}.

\section{What further contact should be proved?}

\begin{exercise}[Every pure entangled two-qubit state violates CHSH]
Put a pure two-qubit state in Schmidt form and optimize planar Pauli
measurements to obtain a CHSH value strictly above two whenever both Schmidt
coefficients are nonzero.
\end{exercise}

\begin{exercise}[Grothendieck turns operators into vectors]
For a real coefficient matrix $(c_{ij})$, formulate the classical sign
optimization and the quantum vector relaxation.  Show how Tsirelson's vector
representation places their ratio under the real Grothendieck constant.
\end{exercise}


\chapter{Infinite observables lead to subfactor index}
\label{ch:subfactors}

\physicalquestion{How can one measure the relative size of nested infinite
observable algebras when ordinary vector-space dimension is meaningless?}

\modelassumptions{Observable algebras act nondegenerately on Hilbert space and
are closed in the weak operator topology.  For the index calculation, the
algebras are finite factors with faithful normalized trace.}

\section{Which closure contains physical limits?}

A von Neumann algebra is a unital $*$-subalgebra
$M\subseteq\mathcal B(H)$ equal to its bicommutant $M''$.  Equivalently, it
is weak-operator closed.  Its predual selects the normal functionals, the
states compatible with monotone limits of observables.

For an inclusion $N\subseteq M$ of finite factors, a trace-preserving
conditional expectation $E_N:M\to N$ is the noncommutative analogue of
averaging onto the smaller observable algebra.  On $L^2(M)$ let $e_N$ be the
orthogonal projection onto $L^2(N)$.

\section{Why do nested expectations produce Temperley--Lieb relations?}

\begin{theorem}[Jones projections]
In the iterated basic construction for a finite-index inclusion, the Jones
projections $e_i$ satisfy
\[
e_i^2=e_i=e_i^*,\qquad e_ie_j=e_je_i\ (|i-j|\ge2),
\qquad e_ie_{i\pm1}e_i=\delta^{-2}e_i,
\]
where $\delta^2=[M:N]$ is the Jones index.
\end{theorem}

\begin{proof}
Idempotence and self-adjointness hold because each $e_i$ is an orthogonal
projection.  Projections separated by at least two stages act on commuting
levels of the tower.  For adjacent stages, use
$e_Nxe_N=E_N(x)e_N$ and the Markov normalization of the trace in the basic
construction; applying the expectation twice contributes the scalar
$[M:N]^{-1}=\delta^{-2}$.  Density of the algebraic tower extends the
relations to $L^2$.
\end{proof}

\section{How does index return to quantum structure?}

The parameter $\delta$ behaves like the dimension of a bimodule and composes
multiplicatively.  Positivity of the Temperley--Lieb trace forces the values
below two to be $\delta=2\cos(\pi/n)$, producing the discrete Jones indices
$4\cos^2(\pi/n)$ \cite{jones-index}.  The same relations yield braid-group representations;
what began as relative size of observable algebras becomes data for particle
exchange.  The quantization theorem is deeper than the displayed relations
and is used here as the bridge to Part VI.

\section{What further contact should be proved?}

\begin{exercise}[Conditional expectation gives a Jones projection]
Prove $e_Nxe_N=E_N(x)e_N$ first on the dense subspace $M\subseteq L^2(M)$.
\end{exercise}

\begin{exercise}[Part V synthesis]
Track one composite-system question through four levels: ordinary norm,
matrix norm, completely positive order, and von Neumann closure.  For each
level identify a physical operation it controls and a counterexample showing
why the preceding level is insufficient.  End by deriving the braid
generators from the Temperley--Lieb projections.
\end{exercise}

\parttransition{Operator spaces and algebras encode states, maps, and
inclusions, but they do not make fusion, reassociation, exchange, defects, and
gluing into first-class operations.  Those processes demand monoidal and
braided categories.}
